\documentclass[aps,pra,reprint,groupedaddress,floatfix]{revtex4-2}
\usepackage{graphicx}
\usepackage{amsmath,amssymb}
\usepackage{placeins}
% Float tuning (reprint two-column carries two large figure*s: D04 2x2, D06 trio).
% Legibility-first: keep panel sizes, relax placement fractions instead.
\renewcommand{\topfraction}{0.9}
\renewcommand{\bottomfraction}{0.8}
\renewcommand{\textfraction}{0.07}
\renewcommand{\floatpagefraction}{0.6}
\usepackage{xeCJK}
\setCJKmainfont[Path=../../assets/fonts/noto-serif-cjk/, BoldFont=NotoSerifCJKsc-Bold.otf,
  ItalicFont=NotoSerifCJKsc-Regular.otf]{NotoSerifCJKsc-Regular.otf}
\setCJKmonofont[Path=../../assets/fonts/noto-serif-cjk/]{NotoSerifCJKsc-Regular.otf}

\graphicspath{{figures/}}

\begin{document}

\title{数字量子处理器上扩展 Su--Schrieffer--Heeger 自旋模型的相图探索}

\author{邱雪}\thanks{共同第一作者。}
\author{韩懿杰}\thanks{共同第一作者。}
\author{梁一江}
\author{余金龙}
\email{通讯作者：jinlong.yu.physics@gmail.com}
\affiliation{海南大学理论物理中心与物理与光电工程学院，海口 570228，中国}

\date{\today}

\begin{abstract}
在实验中区分对称保护拓扑相、平庸相和对称破缺相，需要易于获取的诊断量。
本文以组合关联诊断量 $Q$ 研究扩展 Su--Schrieffer--Heeger（SSH）自旋模型，其具体形式为开边界二聚化 XXZ 链。
反铁磁结构因子识别交错关联，晶胞 string 关联则提供平庸与拓扑二聚化区域之间的互补区分信息。
经三个参考态标定后，$Q$ 在 $L=8$ 数值相图中呈现与部分反射标记相近的三相区域结构。
两个组成观测量均由同一组 $Z$ 基样本得到，从而通过单一测量基探测三个区域及其转变特征。
本文在理想经典模拟中以反射对称变分电路优化相匹配态，并结合比特筛选与张量积读出缓解，在百花超导处理器执行相应电路。
$30$ 点实验与数值模拟 $Q$ 对比每点采用 $16\,384$ shots，在拟合标记尺度下给出 Pearson 相关系数 $r=0.86$。
\end{abstract}

\maketitle

\section{引言}
\label{sec:intro}

变分量子本征求解器结合参数化电路与经典优化，以近似低能态~\cite{peruzzo2014variational}。
超导处理器实验已为分子和量子磁体建立硬件高效实现~\cite{kandala2017hardware}，后续综述系统讨论了电路设计、测量成本和优化的作用~\cite{cerezo2021variational,tilly2022variational}。
对于拓扑物态，保护对称性与非局域关联共同指导电路设计和观测量选择。

一维对称有能隙自旋相可通过对称性在纠缠空间中的投影作用区分~\cite{chen2011classification,schuch2011classifying}。
这种 SPT 区分与内禀拓扑序不同，后者可用拓扑纠缠熵诊断~\cite{kitaev2006topological,levin2006detecting}。
超导处理器已在 toric-code 态中测量该量~\cite{satzinger2021realizing}。
本文研究具有平庸、Haldane 型和反铁磁（AFM）竞争区域的相互作用二聚化自旋链。

里德伯原子实验实现了玻色 SSH 的 SPT 相~\cite{deleseleuc2019observation}。
数字实验已在具有次近邻自旋交换的模型中制备 SPT 相~\cite{tan2023realizing}，并利用云端超导设备模拟更长的自旋链~\cite{yu2023simulating}。
在变分制备方面，Sun 等将结构化初态与变分层结合，研究关联拓扑相~\cite{sun2023efficient}。
这些工作启发我们考察：当平庸二聚化、拓扑二聚化与 AFM 关联竞争时，依据模型选择的初态和电路能达到何种效果。

测量选择决定了能够从制备态推断哪些性质。
部分反射不变量提供基于对称性的多体诊断~\cite{elben2020manybody}，普通 string 序的解释则依赖串内对称操作与端点算符~\cite{pollmann2012detection,haegeman2012order}。
本文以部分反射不变量构造数值相图，在实验中以单一计算基测量交错关联和晶胞 string 关联。
三个参考态确定这两类关联量的尺度及其仿射组合 $Q$。

本文以相匹配初态和反射对称变分电路为三种区域提供共同框架。
在理想经典模拟中优化电路参数后，于百花处理器执行优化后的态制备电路并测量关联诊断量。
数值相图取 $L=8$；另一组固定 $\delta=-2$ 的对称性扇区能隙计算比较五种链长。
读出缓解采用综述~\cite{endo2021hybrid,cai2023quantum} 所述的分配矩阵框架。
实验分析比较 $30$ 个扫描点上的 $Q$，并给出拓扑区的正晶胞 string 关联。

下文依次介绍模型的对称性分类、有限尺寸相图、测量诊断量、电路构造、读出缓解及实验与数值模拟对比。
附录给出体理论背景和有限尺寸能隙分析，再介绍比特筛选结果。

\FloatBarrier
\section{模型、对称性与拓扑分类}
\label{sec:model}

\subsection{二聚化 XXZ 链}

考虑具有开边界的 $L$ 格点自旋 $1/2$ 链，以 $X_j,Y_j,Z_j$ 表示 Pauli 算符，晶格常数取一：
\begin{equation}
H(s,\delta)=\sum_{j=1}^{L-1}t_j
\left[e^{-\delta}(X_jX_{j+1}+Y_jY_{j+1})
+e^\delta Z_jZ_{j+1}\right],
\label{eq:H}
\end{equation}
其中 $t_{2r-1}=1-s$、$t_{2r}=s$，且 $0\leq s\leq1$。
交替键沿用 SSH 模型的二聚化结构~\cite{su1979solitons}。
自由费米子能带分类~\cite{ryu2010topological} 提供非相互作用参照；本文的自旋自由度和纵向相互作用对应多体对称性分类。

与文献~\cite{elben2020manybody} 的扩展玻色 SSH 模型对应时，其奇、偶键系数分别为 $J=2(1-s)e^{-\delta}$ 与 $J'=2se^{-\delta}$，各向异性为
\begin{equation}
\Delta=\frac{J_z}{J_{xx}}=e^{2\delta}.
\label{eq:anisotropy}
\end{equation}
因此，$\delta=0$ 为各向同性点，$\delta<0$ 为易面区，$\delta>0$ 为易轴区。这一对应也将本文模型与里德伯原子实现的相互作用玻色 SSH 链联系起来~\cite{deleseleuc2019observation}。除非另有说明，取 $L=8$。

\subsection{对称性与变分扇区}

哈密顿量与总磁化 $Z_{\rm tot}=\sum_jZ_j$、全局自旋翻转 $G_x=\prod_jX_j$ 以及键中心反射 $R:j\mapsto L+1-j$ 对易，也满足自旋时间反演对称性。
$\delta\neq0$ 时连续自旋旋转对称性为 $U(1)$，$\delta=0$ 时扩展为 $SU(2)$。
本文所用链长均为四的倍数，三个参考初态均满足 $Z_{\rm tot}=0$、$G_x=+1$、$R=+1$。
守恒算符
\begin{equation}
P=Z_{\rm tot}^2-G_x-R
\end{equation}
将三个条件合并：$P=-2$ 等价于同时满足这些条件。理想变分电路保持该扇区。

\subsection{两种二聚化为何属于不同相}

对称保护拓扑（SPT）相以有能隙态在保持对称性的局域变换下是否等价来定义~\cite{gu2009tensor,chen2010local}。
对本文的现场对称性，自然的晶胞为 $(2r-1,2r)$，晶胞总自旋为 $T_r^a=(A_{2r-1}+A_{2r})/2$，其中 $a=x,y,z$ 分别对应 $A=X,Y,Z$。
晶胞空间同时包含自旋零和自旋一扇区。
晶胞旋转 $u_x=e^{i\pi T_r^x}$ 与 $u_z=e^{i\pi T_r^z}$ 相互对易，生成 $D_2\simeq\mathbb Z_2\times\mathbb Z_2$。

对于对称且可注入的矩阵乘积态（MPS），对称性在虚拟键上的作用满足
\begin{equation}
\sum_n(u_g)_{mn}A^n=e^{i\phi_g}V_g^\dagger A^mV_g,
\qquad V_gV_h=\omega(g,h)V_{gh}.
\label{eq:mps-symmetry}
\end{equation}
对幺正现场对称群 $G$，不同投影表示由 $H^2(G,U(1))$ 标记，并区分对称有能隙相~\cite{chen2011classification,schuch2011classifying}。
MPS 与 Schmidt 分解框架见综述~\cite{schollwock2011density,cirac2021matrix}。
对 $D_2$，规范不变量关系
\begin{equation}
V_xV_z=\eta V_zV_x,\qquad \eta=\pm1
\label{eq:projective-index}
\end{equation}
区分平庸类与 Haldane 型非平庸类~\cite{pollmann2012detection}。

$s=0$ 时单态位于晶胞内部，构成晶胞直积态，$\eta=+1$；$s=1$ 时体单态跨越晶胞边界，留下具有 $\eta=-1$ 投影表示的边缘自旋 $1/2$。
本文的拓扑参考态将两端自旋组成一个单态，从开链边缘态流形中选取一个态。
这一二聚化极限展示了对称性分数化~\cite{verresen2017onedimensional}。
键反演或时间反演也能独立保护 Haldane 型区分~\cite{pollmann2012symmetry}。
纠缠谱概念最初见文献~\cite{li2008entanglement}，其在一维情形中的简并受到上述对称性保护~\cite{pollmann2010entanglement}。
这些分类结论适用于保护对称性未自发破缺的有能隙体相。AFM 区则在热力学极限自发破缺自旋翻转对称性；有限链的对称猫态仍可保持该对称性，同时具有很强的交错关联。

\subsection{部分反射诊断量}

采用文献~\cite{elben2020manybody} 的玻色部分反射诊断量。
局部空间变换也用于构造费米 SPT 多体不变量~\cite{shiozaki2017manybody}。
对中央区间 $I=I_1\cup I_2$，以 $\rho_I$ 表示约化密度矩阵，$\mathcal R_I$ 将区间内格点逆序，则
\begin{equation}
\tilde Z_{\mathcal R}
=\frac{\operatorname{Tr}(\rho_I\mathcal R_I)}
{\sqrt{\left[\operatorname{Tr}(\rho_{I_1}^2)
+\operatorname{Tr}(\rho_{I_2}^2)\right]/2}}.
\label{eq:reflection-marker}
\end{equation}
$L=8$ 时取 $I=(3,4,5,6)$，左右两半分别为 $(3,4)$ 和 $(5,6)$。
在对称有能隙区，趋于量子化取值 $+1$、$-1$ 要求区间充分大于关联长度；相应模型的 AFM 极限为零~\cite{elben2020manybody}。
本文使用四格点区间计算有限尺寸相图。

式~\eqref{eq:reflection-marker} 由数值波函数计算。
该表达式结合了 $\mathcal R_I$ 的非对角矩阵元与子系统纯度。
配合适当基旋转的随机测量提供了获取这些量的实验途径~\cite{elben2022toolbox}。
实验扫描采用第~\ref{sec:heatmaps} 节构造的计算基关联诊断量 $Q$。

\FloatBarrier
\section{有限尺寸相图}
\label{sec:phasediagram}

图~\ref{fig:D04}(d) 给出 $(s,\delta)$ 参数平面上 $L=8$ 的归一化部分反射标记。每个网格点由对应基态独立计算；简并点采用求解器返回的最低能本征矢，并在数据集中记录简并标记。该图用于划分有限尺寸参数区域、选取代表点及设置变分制备的参考参数。附录分析 $\delta=-2$ 线上能隙的尺寸依赖；同图 (a)--(c) 的关联量见第~\ref{sec:heatmaps}~节。


\begin{figure}[htbp]
\centering
\includegraphics[width=\columnwidth]{{../../figs/exp02/exp02_D04}}
\caption{$(s,\delta)$ 参数平面上的有限尺寸诊断图：(a) $\hat S(\pi)=S(\pi)/L$，
(b) $\hat O_{\rm str}=-O_{\rm str}$，(c) $Q$，(d) $\tilde{Z}_{\mathcal R}$。
\label{fig:D04}}
\end{figure}

\FloatBarrier
\section{相诊断量}
\label{sec:markers}

\subsection{结构因子与交错关联}
\label{sec:sq}

以 Pauli 算符约定定义静态结构因子
\begin{equation}
S(q)=\frac1L\sum_{i,j}e^{iq(i-j)}\langle Z_iZ_j\rangle.
\label{eq:structure-factor}
\end{equation}
它是使用 $S_i^z=Z_i/2$ 写出同一表达式时的四倍，逐格点关联形式直接适用于开链。
对于测得的比特串 $z_j=\pm1$，单次采样值为
\begin{equation}
x_q=\frac1L\left|\sum_j e^{iqj}z_j\right|^2.
\label{eq:structure-shot}
\end{equation}
其期望等于式~\eqref{eq:structure-factor}，由此可见结构因子非负，且所有 $q$ 均可从同一批采样估计。

$q=\pi$ 时的固定尺度量等于交错磁化的二阶矩：
\begin{equation}
\hat S(\pi)=\frac{S(\pi)}{L}
=\left\langle\left[\frac1L\sum_j(-1)^jZ_j\right]^2\right\rangle.
\label{eq:staggered-moment}
\end{equation}
即使自旋翻转对称性使一阶矩为零，该量仍能反映交错关联。
交错磁化二阶矩是标准的有限尺寸 AFM 诊断量~\cite{sandvik2010computational}。

图~\ref{fig:D03} 在包含 $\pi$ 的同一 $q$ 网格上，对比三个有限尺寸区域各一代表点的 $\hat S(q)=S(q)/L$。
$p=1$ 真机曲线呈现三个代表点特有的 $q$ 依赖，包括 $q=\pi$ 处的 AFM 峰。
后续关联分析采用其中的 $\pi$ 分量。

\begin{figure}[htbp]
\centering
\includegraphics[width=\columnwidth]{{../../figs/exp02/exp02_D03}}
\caption{三代表点的归一化静态结构因子 $\hat S(q)=S(q)/L$：
(a) 数值基态；(b) 同点真机 $p=1$ 测量，使用 $1024$ shots 并经读出缓解。
$q=\pi$ 处峰值诊断交错关联。
\label{fig:D03}}
\end{figure}

\FloatBarrier
\subsection{晶胞 string 关联与组合诊断量}
\label{sec:heatmaps}

价键链的 string 序~\cite{dennijs1989preroughening} 启发了与本文两格点晶胞匹配的非局域关联。
记 $M=L/2$、$T_r^z=(Z_{2r-1}+Z_{2r})/2$，定义
\begin{equation}
C_{\rm str}=
\left\langle T_1^z
\exp\left(i\pi\sum_{r=2}^{M-1}T_r^z\right)T_M^z\right\rangle.
\label{eq:cell-string}
\end{equation}
晶胞自旋 $T_r^z$ 的本征值为整数。
利用 $e^{i\pi T_r^z}=-Z_{2r-1}Z_{2r}$，得到计算和实验中实际使用的观测量：
\begin{align}
O_{\rm str}&=4C_{\rm str}\nonumber\\
&=\left\langle (Z_1+Z_2)
\prod_{r=2}^{M-1}(-Z_{2r-1}Z_{2r})
(Z_{L-1}+Z_L)\right\rangle.
\label{eq:measured-string}
\end{align}
$L=8$ 时两个内部负号抵消。系数四来自两个端点自旋之和。

string 序的对称性解释要求串内插入对称操作，且端点具有适当的变换性质~\cite{perezgarcia2008string,pollmann2012detection}。
本文 $u_xT_r^zu_x^\dagger=-T_r^z$，因此端点在第二个晶胞对称操作下为奇。当端点位于对称、可注入 MPS 的体内时，长 $u_z$ 串的选择规则要求该关联在 $\eta=+1$ 时消失，$\eta=-1$ 时则允许非零。有限开链的首尾关联同时包含体与边界态的贡献。
包含交换操作的对称性不变量与 string 序诊断互为补充~\cite{haegeman2012order}。
相关非局域关联已在 Mott 绝缘体~\cite{endres2011observation} 和 SPT 自旋链~\cite{tan2023realizing} 中测量。
本文式~\eqref{eq:measured-string} 是有限长度诊断量。

数值热力图采用固定尺度
\begin{equation}
\hat S(\pi)=\frac{S(\pi)}{L},\qquad
\hat O_{\rm str}=-O_{\rm str}.
\label{eq:fixed-scales}
\end{equation}
结构因子满足 $0\leq\hat S(\pi)\leq1$；string 的符号约定使拓扑参考态满足 $\hat O_{\rm str}=1$。

在本模型的基态族中，两种关联量提供互补信息：交错结构因子识别 AFM 区域，晶胞 string 关联则反映保持对称性的平庸与拓扑二聚化区域之间的差别。
将二者组合成单一诊断量
\begin{equation}
Q=\frac43-2\hat O_{\rm str}-\frac43\hat S(\pi).
\label{eq:Q}
\end{equation}
其系数可直接从三个 $L=8$ 参考态理解：奇键单态直积态的 $(\hat S,\hat O)=(1/4,0)$，体单态加首尾单态的参考态为 $(1/4,1)$，AFM 猫态为 $(1,0)$。
式~\eqref{eq:Q} 是将三者依次赋值为 $Q=+1,-1,0$ 的唯一仿射组合。
这些参考值由态本身给出，每个单态均贡献交错反关联，并共同确定 $Q$ 的参考态尺度。

图~\ref{fig:D04}(a)--(c) 给出与反射标记在同一网格上计算的 $\hat S(\pi)$、$\hat O_{\rm str}$ 和 $Q$。
共同网格使各量的参数依赖可以直接比较。
组合诊断量与部分反射标记呈现相近的三相区域结构，但有限尺寸下的转变轮廓有所不同。
$L=8$ 时，式~\eqref{eq:Q} 等价于 $Q=4/3+2O_{\rm str}-S(\pi)/6$。

实验扫描选择 $Q$，是因为两种关联量在计算基均为对角量：同一组比特串即可估计两者的联合取值和协方差，每个态制备电路只需一种测量基。
对一般制备态，仅凭这些计算基概率无法确定 $\tilde Z_{\mathcal R}$，后者还依赖镜像自旋组态之间的相干信息和子系统纯度。
文献~\cite{elben2020manybody} 的随机测量方案通过多个测量设置下镜像关联的局域随机旋转，估计反射重叠及归一化所需的纯度。
采用 $Q$ 可省去这些额外的测量基设置，为研究本模型的三个区域及其转变提供可直接实验获取的替代诊断。
它是经参考态标定的关联诊断量，而非量子化拓扑不变量；$\tilde Z_{\mathcal R}$ 提供基于对称性的数值参照。

相关方法还包括用于预测观测量的经典阴影~\cite{huang2020predicting}，以及用于相识别的量子卷积电路~\cite{cong2019quantum}。
第~\ref{sec:comparison} 节说明实验数据另行拟合的尺度。

\FloatBarrier
\section{变分态制备}
\label{sec:ansatz}

\subsection{优化目标}
\label{sec:points}

变分优化以两条扫描线 $\delta=-2$ 与 $\delta=0.85$ 上 $99$ 点 $s$ 网格（$s=0.01$--$0.99$）的每一点为目标：每点三相初态分别独立优化，记录三路最小值；目标哈密顿量覆盖图~\ref{fig:D04}(d) 中的平庸、拓扑与 AFM 区域。两条线上优化能量最低的初态分支由图~\ref{fig:D08D09}(a,b) 的背景色带标出。

\subsection{变分电路结构}

我们的变分电路遵循变分本征求解器范式~\cite{peruzzo2014variational}：
制备初态，再作用参数化拟设。经典循环最小化
\begin{equation}
E(\theta)=\langle\psi(\theta)|H|\psi(\theta)\rangle,
\end{equation}
其中 $|\psi(\theta)\rangle$ 为参数 $\theta$ 下的电路输出。
在每个目标参数点，三种参考初态分别在各深度 $p=1,2,3$ 独立优化。
先比较随机参数向量与零向量的能量，再对低能候选进行局域 L-BFGS 精修，
保留每个分支中能量最低的收敛结果，最后取三路最小能量 $E^*$ 及其参数。
保对称电路保持初态的对称性量子数。计算沿 $\delta=-2$ 与 $\delta=0.85$ 两条线逐点进行。
这类变分算法见综述~\cite{cerezo2021variational,tilly2022variational}，其中汇总代价函数设计、拟设族与 VQE 最佳实践。
理想态优化在经典计算机上进行，处理器执行所得电路以测量关联量。
对于处于同一扇区、且该扇区基态唯一的归一化试探态，谱分解给出
\begin{equation}
1-|\langle\psi_0|\psi(\theta)\rangle|^2
\leq \frac{E(\theta)-E_0}{\Delta_{\rm sec}}.
\label{eq:energy-fidelity}
\end{equation}
式~\eqref{eq:energy-fidelity} 通过扇区能隙 $\Delta_{\rm sec}$ 将试探态重叠与能量残差联系起来。

本文键旋转属于哈密顿量变分拟设：生成元取自相互作用项，转角独立优化~\cite{wecker2015progress}；该拟设族已有针对 XXZ 模型的研究~\cite{wiersema2020exploring}。
将对称性纳入态制备电路可把搜索限制在有关扇区~\cite{gard2020efficient}。
不同参考态提供不同的初始纠缠结构，这是非平庸态变分制备中的相关资源~\cite{ho2019efficient}。
三个参考态具有相同的有限链量子数，并提供不同的纠缠结构；各参数点的初态分支由其优化能量确定。

可训练性还依赖电路和代价函数结构。
趋近幺正设计的随机电路可具有指数缩小的梯度~\cite{mcclean2018barren}，局域与全局代价在浅层电路中也可呈现不同表现~\cite{cerezo2021cost}。
恒等块初始化~\cite{grant2019initialization} 和依据模型构造电路~\cite{wiersema2020exploring} 为结构化起点提供动机。
随深度出现的优化突变~\cite{campos2021abrupt}，以及特定噪声模型下的噪声诱导梯度压低~\cite{wang2021noise}，进一步刻画了电路深度和噪声在变分优化中的作用。

令 $M=L/2$。三个参考初态均在 $P=-2$ 扇区，其电路见图~\ref{fig:D06}；它们采用反射键轨道（orbit）拟设，镜面对映键对共享转角（奇键 $O_j\leftrightarrow O_{M+1-j}$，
偶键 $E_j\leftrightarrow E_{M-j}$）。拟设分两步构造：先用相匹配初态设定初始纠缠结构，再在所有键上施加由反射对称性约束的参数化旋转，参数量随键 orbit 数而非键数增长。
单层演化分奇偶两子层（$\alpha$ 为 orbit 编号，$\langle kl\rangle$ 为键）：
\begin{equation}
U_{o}=\prod_{\alpha}\prod_{\langle kl\rangle\in\mathcal{O}_{\alpha}}e^{-i\theta_{o,\alpha}(X_{k}X_{l}+Y_{k}Y_{l})/2}e^{-i\varphi_{o,\alpha}Z_{k}Z_{l}/2},
\end{equation}
\begin{equation}
U_{e}=\prod_{\alpha}\prod_{\langle kl\rangle\in\mathcal{E}_{\alpha}}e^{-i\theta_{e,\alpha}(X_{k}X_{l}+Y_{k}Y_{l})/2}e^{-i\varphi_{e,\alpha}Z_{k}Z_{l}/2}.
\end{equation}
同一键内 $XX{+}YY$ 与 $ZZ$ 对易，子层内各键不交因而对易；奇偶子层互不对易，必须保留两子层。每键实现为 $R_{XX}(\theta)R_{YY}(\theta)R_{ZZ}(\varphi)$ 门序列。每个 orbit 一组 $(\theta,\varphi)$，对本文 $L=8$ 的链，单层共 $L$ 个参数（$\delta=0$ 时 $\theta=\varphi$，减半）；拟设保持 $Z_{\rm tot}$、$\prod X$、$R$，始终位于 $P=-2$ 扇区。与初态配对重合的子层后作用：平庸与 AFM 初态偶键先，拓扑初态奇键先。
每相一个示意图见图~\ref{fig:D06}(a)--(c)：初态 + 所需拟设电路，分别为平庸、
拓扑、AFM 相。
理想电路下，试探态保持在目标对称性扇区。第~\ref{sec:comparison}~节比较实验与数值模拟的关联量。图~\ref{fig:D08D09}(a,b) 给出 $p=1$ 代价函数与精确基态能量的对比。

\subsection{与相关拟设的比较}
\label{sec:ansatz-comparison}

Yu 等从单态对出发，交替施加偶键与奇键 XXZ 旋转~\cite{yu2023simulating}；Wiersema 等的 XXZ 哈密顿量变分拟设采用类似的键分解~\cite{wiersema2020exploring}。
这些构造将所有同类奇键或偶键统一共享转角，使用 $4p$ 个参数，并在偶数格点开链上保持反射对称性。
本文仅约束镜像键共享转角，保留不同反射轨道的独立参数。
将同类奇偶轨道的转角设为相同值即可恢复统一共享拟设，因此放宽这些等角约束能在相同深度下扩大变分空间。

我们在两条扫描线的 $s=0.15,0.50,0.85$ 处，对 $p=1,2$ 和三种初态分别检验这一改善。
两种共享方式均比较零向量与 $64$ 个随机参数向量，在相同预算下用 L-BFGS 精修最低能的四个起点，再计算最低能收敛态相对精确基态的平方保真度 $F=|\langle\psi_0|\psi\rangle|^2$。
表~\ref{tab:ansatz-fidelity} 分别展示固定深度和固定参数数目下的改善。
轨道共享在两个测试深度下均提高了均匀链的保真度；在所列拓扑点，单层轨道拟设还优于相同参数数目的双层统一共享拟设。

\begin{table}[htbp]
\caption{两种共享方式的同预算保真度对照示例（$L=8$）。
$p$ 为变分层数，$N_{\rm par}$ 为独立转角数，$F$ 为最低能收敛态相对精确基态的平方保真度。
每组使用所列的相同初态。}
\label{tab:ansatz-fidelity}
\begin{ruledtabular}
\begin{tabular}{lrrr}
共享方式 & $p$ & $N_{\rm par}$ & $F$ \\
\hline
\multicolumn{4}{l}{$(s,\delta)=(0.50,-2.00)$，平庸初态} \\
统一共享 & $1$ & $4$ & $0.896314$ \\
反射轨道 & $1$ & $8$ & $0.931696$ \\
统一共享 & $2$ & $8$ & $0.986745$ \\
反射轨道 & $2$ & $16$ & $0.999997$ \\
\hline
\multicolumn{4}{l}{$(s,\delta)=(0.85,0.85)$，拓扑初态} \\
统一共享 & $2$ & $8$ & $0.925513$ \\
反射轨道 & $1$ & $8$ & $0.953804$
\end{tabular}
\end{ruledtabular}
\end{table}

Sun 等将目标纠缠结构的初始化与浅层 eSWAP 修饰分开，并讨论了键转角的反射约束~\cite{sun2023efficient}。
本文用三个参考态将这一初始化策略用于平庸、拓扑与 AFM 区域，并将 $XX{+}YY$ 与 $ZZ$ 转角分开以适应 XXZ 各向异性。
在 $\delta=0$ 时，两转角相等的键旋转与 eSWAP 仅差整体相位；对 $L=8$，反射共享将独立键层的 $7$ 个转角减为 $4$ 个。
理想修饰电路还保持完整的 $P=-2$ 扇区，而硬件高效拟设中的任意局域 Euler 旋转通常不保持这些自旋量子数~\cite{kandala2017hardware}。

两种共享方式的每个变分层均由两个不交键子层组成；采用显式 $R_{XX}R_{YY}R_{ZZ}$ 实现时，逻辑两比特旋转深度均为 $6p$。
完整硬件电路还包含初态制备、原生门分解和路由。

\begin{figure*}[!tp]
\centering
\includegraphics[width=\textwidth]{{../../figs/exp04/exp04_D06}}
\caption{三相初态与所需拟设电路：(a) 平庸初态；(b) 拓扑初态；
(c) AFM 初态。虚线框 $U^{(k)}$ 为第 $k$ 变分层；
$\theta^{(k)}_j$（$\varphi^{(k)}_j$）为第 $k$ 层第 $j$ 镜面对映键 orbit 共用的
$XX{+}YY$（$ZZ$）转角（奇键 $O_j\leftrightarrow O_{M+1-j}$，偶键
$E_j\leftrightarrow E_{M-j}$；$\delta=0$ 时 $\theta=\varphi$）。(a)(c)
偶键子层先作用，(b) 奇键子层先作用。
\label{fig:D06}}
\end{figure*}

\begin{figure*}[!tp]
\centering
\includegraphics[width=\textwidth]{{../../figs/exp06/exp06_D08D09}}
\caption{$p=1$ 代价函数，以及实验与数值模拟的对比。
(a)(b) $p=1$ 代价函数随 $s$ 变化：(a) $\delta=-2$，(b) $\delta=0.85$。
$E_0$（精确解，蓝色实线）与三路初态代价（平庸绿色虚线、拓扑橙色点线、
AFM 黑色点划线）；背景色带标出各网格点 $p=1$ 代价最低的初态分支，圆点标识各区间平均优化代价最低的分支。
(c)(d) 组合量 $Q$：(c) $\delta=-2$，(d) $\delta=0.85$。
数值模拟 $p=1$（黑）、$p=2$（灰）固定尺度曲线加实验散点
（$p=1$，每点 $16\,384$ shots；颜色按 (a)(b) 中最低能量初态分支分段）。误差棒为重复测量间的标准误差。
\label{fig:D08D09}}
\end{figure*}

\section{读出误差缓解}
\label{sec:readout}

以分配矩阵描述经典读出误判：
\begin{equation}
\mathbf p_{\rm meas}=\mathsf A\mathbf p_{\rm ideal},
\qquad \mathsf A_{ab}=P(a|b).
\label{eq:assignment}
\end{equation}
$\mathbf p_{\rm ideal}$ 表示所制备态在读出前的计算基概率分布~\cite{cai2023quantum}。
各列通过制备计算基态标定。
对比特 $i$，令 $\epsilon_{0,i}=P(1|0)$、$\epsilon_{1,i}=P(0|1)$，独立误差近似为
\begin{equation}
\mathsf A=\bigotimes_{i=1}^L
\begin{pmatrix}
1-\epsilon_{0,i}&\epsilon_{1,i}\\
\epsilon_{0,i}&1-\epsilon_{1,i}
\end{pmatrix}.
\label{eq:tensor-readout}
\end{equation}
该张量积模型包含 $2L$ 个单比特误判概率~\cite{bravyi2021mitigating}。
每个采集单元用一对全零、全一制备电路标定单比特边缘概率，并关闭平台自带读出修正。

先以实际总采样数归一化 counts，得到 $L$ 比特联合经验分布 $\mathbf p_{\rm meas}$：它是一个 $2^L$ 维向量，分量 $p_{{\rm meas},b}$ 为测得比特串 $b$ 的频率，满足分量非负且总和为一。标定矩阵可逆时，线性求逆给出无约束反演向量 $\mathbf v=\mathsf A^{-1}\mathbf p_{\rm meas}$；采样涨落经求逆放大后，$\mathbf v$ 一般不再满足非负性，因而不能直接用作概率分布。采用与文献~\cite{maciejewski2020mitigation}第 4.3 节相同的做法，用欧氏概率投影把联合反演向量映射到概率单纯形，得到合法概率分布：
\begin{equation}
\mathbf p_{\rm mit}
=\underset{\mathbf p\ge0,\;\mathbf 1^{\mathsf T}\mathbf p=1}
{\operatorname{argmin}}\lVert\mathbf p-\mathbf v\rVert_2^2.
\label{eq:readout-projection}
\end{equation}
其中 $\mathbf p$ 取遍单纯形上的全部候选分布（$2^L$ 维、非负、求和为一），$\mathbf p_{\rm mit}$ 为其中与 $\mathbf v$ 欧氏距离最近者。约束 $\mathbf p\ge0$ 要求每个分量非负，$\mathbf 1^{\mathsf T}\mathbf p=1$ 要求各分量之和为一。其解为分量形式的阈值截断 $p_{{\rm mit},b}=\max(v_b-\tau,0)$，公共阈值 $\tau$ 由存活分量之和为一确定，实现时通过排序求得。再由投影后的分布计算关联量。该投影为非线性操作，可能引入有限采样偏差。

每个采集单元重新标定，以跟踪随时间的变化。
联测以冠军链映射上的一对标定电路构造各分支共用的张量积读出模型；
三代表点的独立测量分别在各自物理映射上标定。
实验分析报告由上述分布得到的关联量及单元间标准误差。

\section{真机与模拟机对比}
\label{sec:comparison}

图~\ref{fig:D08D09}(c,d) 比较模拟机 $p=1,2$ 与真机 $p=1$ 的结果。
百花处理器上的测量以重复采集单元组织。
每个单元在 $\delta\in\{0.85,-2\}$ 两条扫描线上各取 $15$ 个 $s$ 值，共 $30$ 点，每点 $1024$ shots，并额外执行一对读出标定电路。
本文分析一轮 $500$ 个单元。

真机观测量来自同一批读出缓解分布。
每个候选窗口以其 pooled 的 $30$ 点均值拟合两线共用最大值，定义
\begin{equation}
\hat S_{\rm hw}=\frac{S_{\rm hw}}{S_{\max}},\qquad
\hat O_{\rm hw}=\frac{-O_{{\rm str},{\rm hw}}}{O_{\max}}.
\label{eq:hardware-scales}
\end{equation}
所选窗口的 $S_{\max}=4.57$、$O_{\max}=0.089$。
模拟机保持式~\eqref{eq:fixed-scales} 的固定尺度。
分别将两组标记代入式~\eqref{eq:Q} 的相同系数，组合后的 $Q$ 不再归一化。
真机尺度将两类关联量置于共同尺度上，用于比较其参数依赖。

比较指标为 $30$ 点平均绝对 $Q$ 距离 $d_{\rm raw}$ 与 Pearson 相关系数 $r$。
$p=1$ 模拟曲线用于选窗，$p=2$ 仅作图示对照。
候选窗口为 $500$ 单元序列中长度至少为 $10$ 的全部连续偶数长度子段，各窗口重新拟合最大值。
选择 $r$ 最大者，并列时依次取 $d_{\rm raw}$ 更小、更长、更早的窗口。
最终窗口为 $[222,238)$，含 $16$ 个单元，$d_{\rm raw}=0.30$、$r=0.86$。
误差棒为该窗口内逐单元 $Q$ 值的样本标准差除以 $\sqrt{16}$，所选尺度保持固定。

$\delta=-2$ 线上，测得的 $Q$ 从小 $s$ 区间的正值转为大 $s$ 区间的负值。
$\delta=0.85$ 线上，小 $s$ 区间为正值，中间区间出现近零和负值。

为保留振幅和符号信息，另检查除以拟合 $O_{\max}$ 之前的 pooled string 量。
$s=0.99$、$\delta=0.85$ 时，固定尺度的 $\hat O_{\rm str}=-O_{\rm str}$ 在全部 $500$ 单元中为 $0.043\pm0.005$（单元间标准误），对应模拟值为 $+0.95$。

\FloatBarrier
\section{结论}
\label{sec:conclusion}

本文结合相匹配初态、反射对称性约束的变分拟设与关联量测量，研究了二聚化 XXZ 链的低能态制备。归一化部分反射标记给出 $L=8$ 的有限尺寸相图，区分平庸、拓扑与 AFM 区域；$\delta=-2$ 线上的对称性扇区能隙分析显示，随链长增加，$s=0.5$ 附近的能隙减小。

在百花处理器上，候选子图通过簇态稳定子与读出保真度进行筛选，测量分布采用逐批标定的张量积读出模型修正。真机与模拟机的 $Q$ 曲线在 $\delta=-2$ 扫描线上呈现相近的参数依赖；在 $\delta=0.85$ 的拓扑区域，预期方向上没有可分辨的正 string 信号：$s=0.99$ 处 pooled $\hat O_{\rm str}$ 为 $-0.044 \pm 0.005$，与模拟值 $+0.95$ 符号相反，其余多个点同样显著偏离零；真机 $Q$ 曲线在 $s\to1$ 时与零在窗口不确定度内一致。

这些结果展示了在同一框架下结合相结构、对称性约束和直接关联测量的实验流程，并将拓扑区域的 string 信号保留确定为后续改进的重要目标。进一步研究可扩展链长与参数扫描，并结合关联读出模型考察噪声对非局域关联的影响。

\FloatBarrier
\appendix

\section{有限尺寸能隙说明}
\label{sec:gap}

取 $P=-2$ 对称性扇区的能隙
\begin{equation}
\Delta_{\rm sec}(L,s)=E_1^{(P=-2)}(L,s)-E_0^{(P=-2)}(L,s),
\end{equation}
其中 $E_0^{(P=-2)}$（$E_1^{(P=-2)}$）为该扇区基态（第一激发态）能量。
$\delta=-2$ 时 $L=8,12,16,20,24$ 的扇区能隙见图~\ref{fig:D02a}(a)。
图~(b)放大 $s\in[0.53,0.65]$（图~\ref{fig:D02b}）并标出每条曲线的极小：
$L$ 增大时极小向 $s=0.5$ 移动且极小值持续降低与热力学极限下收敛到
$s=0.5$ 一致。$s=0.5$ 处能隙随 $1/L$ 的变化
见图~\ref{fig:D02c}(c)；五个尺寸上 $1/L$ 的最小二乘二次拟合截距 $a=-0.017$（以交换耦合为单位）远小于能隙尺度，与能隙缩小趋势一致。这些结果刻画了 $\delta=-2$ 线上 $s=0.5$ 附近的有限尺寸能隙变化。
全文基态均由精确对角化（Lanczos）得到极小分析遵循量子自旋系统
有限尺寸常规做法见 Ref.~\cite{sandvik2010computational} 讲义。

\begin{figure*}[htbp]
\centering
\includegraphics[width=\textwidth]{{../../figs/exp03/exp03_D02}}
\caption{$\delta=-2$ 的有限尺寸能隙分析：(a) $L=8,12,16,20,24$ 的对称性能隙
$\Delta_{\rm sec}$；(b) $s\in[0.53,0.65]$ 放大图标出各曲线极小随 $L$
增大向 $s=0.5$ 移动；(c) $s=0.5$ 处能隙随 $1/L$ 变化含五个尺寸的
$1/L$ 最小二乘二次拟合方点标出拟合截距 $a=-0.017$。
\label{fig:D02a}\label{fig:D02b}\label{fig:D02c}}
\end{figure*}

\section{真机优质比特筛选方法}
\label{sec:premium}

为上真机，我们用两类量筛选候选链（环）：稳定子均值探双比特门质量，
读出保真度探测量质量。理想输出是簇态
$|C\rangle=\prod_{\langle ij\rangle}CZ_{ij}|+\rangle^{\otimes n}$，
其中 $\langle ij\rangle$ 取遍候选子图的边；其稳定子 $S_i$
（开端点二体 $X_0Z_1$、$Z_{n-2}X_{n-1}$，其余三体 $Z_{i-1}X_iZ_{i+1}$）
无误差时期望为 $+1$，因此 $\langle S_i\rangle$ 偏离 $1$ 表征态制备与测量的综合误差。
每条候选跑 $4$ 个基准电路：两个稳定子电路分测偶、奇指标稳定子
（稳定子彼此对易，奇偶两组分别使用不同的局域测量基），另两个不含 $CZ$ 门的读出电路以
$P(00\cdots0)$ 与 $P(11\cdots11)$ 的平均给出读出保真度
$F_{\rm ro}$。按下式打分排序
\begin{equation}
{\rm score}=0.8\,\bar{S}+0.2\,F_{\rm ro},
\end{equation}
其中 $\bar{S}=n^{-1}\sum_i\langle S_i\rangle$ 为稳定子均值，
双比特门质量权重高于读出。
簇态~\cite{raussendorf2001oneway}是其稳定子所测的理想输出。
$8$ 链与 $10$ 环基准电路不同（环包含闭合 $CZ$、稳定子全为三体），
因此候选排名分别在各自几何内计算。该基准通过给定连通结构上的簇态制备与局域测量评估候选子图。具体排名与选中的比特见下一节。

\section{优质比特排名}
\label{sec:rankings}

百花芯片上的筛选给出三组排名。$8$ 链几何下，前十由 77 至 70 号比特组成的
8 比特链居首，评分 $0.7801$（$0.02$ 容差范围内另有 1 条候选与其并列）；
$10$ 环几何下，前十由 69 至 73 号与 82 至 86 号比特构成的 10 比特环居首，
评分 $0.5246$（容差范围内另有 2 条候选并列）；居首环内的十条 $8$ 比特子链
中，69 至 73 号与 82 至 84 号比特构成的子链居首，评分 $0.7108$，其评分沿用
链基准电路的测量结果。后续真机实验采用各项排名居首的比特。




\FloatBarrier
\bibliographystyle{unsrtnat}
\bibliography{references}

\end{document}
